% Partial boundary observation: relative homology kernel, minimal audit completion, and matroidal structure.

\paragraph{部分边界观测的核与审计补全}
在定理 \ref{thm:spg-boundary-godelization-holographic-dictionary} 的“全边界可见”情形下，顶维边界单射给出严格可逆的离散全息字典。
当仅能观测部分取向面（屏幕）时，信息缺口由一个相对同调群精确刻画，并且其秩同时控制最小审计补全成本与可达的 Gödel 型补账寄存器维数。

\begin{theorem}[部分边界观测的屏幕核与相对顶维同调]\label{thm:spg-partial-boundary-screen-kernel-relative-homology}
固定 \(m,n\ge 1\) 并沿用定理 \ref{thm:spg-dyadic-cubical-boundary-injective} 的 dyadic 立方复形 \(\mathcal Q_m\) 与链群 \(C_k(\mathcal Q_m;\ZZ)\) 记号。
令 \(\overrightarrow{\mathcal F_{m}^{n-1}}\) 为全部取向 \((n-1)\)-胞腔所成的标准基。
对任意 \(S\subset \overrightarrow{\mathcal F_{m}^{n-1}}\)，令
\[
\Pi_S:C_{n-1}(\mathcal Q_m;\ZZ)\longrightarrow \ZZ^{S}
\]
为对 \(S\) 坐标的投影，并定义部分边界观测算子
\[
f_S:=\Pi_S\circ\partial_n:\ C_n(\mathcal Q_m;\ZZ)\longrightarrow \ZZ^{S}.
\]
令 \(A_{\neg S}\subset \mathcal Q_m^{(n-1)}\) 为从 \((n-1)\)-骨架中删去 \(S\) 中那些 \((n-1)\)-胞腔后得到的子复形（保留全部 \(\le n-2\) 维胞腔，因而仍为子复形，并且 \(C_n(A_{\neg S};\ZZ)=0\)）。
则有自然同构
\[
\ker f_S\ \cong\ H_n(\mathcal Q_m,\ A_{\neg S};\ZZ).
\]
\end{theorem}
\begin{proof}
由于 \(C_n(A_{\neg S};\ZZ)=0\)，相对链群满足 \(C_n(\mathcal Q_m,A_{\neg S};\ZZ)=C_n(\mathcal Q_m;\ZZ)\)。
相对边界算子为
\[
\partial_n^{\mathrm{rel}}:C_n(\mathcal Q_m;\ZZ)\longrightarrow C_{n-1}(\mathcal Q_m;\ZZ)/C_{n-1}(A_{\neg S};\ZZ).
\]
因此 \(x\in C_n(\mathcal Q_m;\ZZ)\) 为相对 \(n\)-循环当且仅当 \(\partial_n x\in C_{n-1}(A_{\neg S};\ZZ)\)，等价于 \(\Pi_S(\partial_n x)=0\)，亦即 \(x\in\ker f_S\)。
从而 \(\ker f_S=Z_n(\mathcal Q_m,A_{\neg S};\ZZ)\)。
\par
另一方面，\(\mathcal Q_m\) 无 \((n+1)\)-胞腔，故 \(C_{n+1}(\mathcal Q_m,A_{\neg S};\ZZ)=0\) 并且 \(B_n(\mathcal Q_m,A_{\neg S};\ZZ)=\mathrm{im}\,\partial_{n+1}^{\mathrm{rel}}=0\)。
因此
\[
H_n(\mathcal Q_m,A_{\neg S};\ZZ)=Z_n(\mathcal Q_m,A_{\neg S};\ZZ)/B_n(\mathcal Q_m,A_{\neg S};\ZZ)\cong \ker f_S.
\]
证毕。
\end{proof}

\begin{corollary}[最小审计补全成本等于屏幕核秩]\label{cor:spg-partial-boundary-min-audit-cost-kernel-rank}
沿用定理 \ref{thm:spg-partial-boundary-screen-kernel-relative-homology} 的记号，并令 \(r:=\rank_{\ZZ}(\ker f_S)\)。
则任意使 \(f_S\) 可单射提升的审计补全 \((R,r_S)\)（定义 \ref{def:cdim-audit-completion-loss}）都满足
\[
\cdim(R)\ \ge\ r,
\]
并且该下界可达。
因此，\(f_S\) 的最小圆维审计补全成本等于
\[
\min_{(R,r_S)}\cdim(R)=\rank_{\ZZ}(\ker f_S)=\rank_{\ZZ}H_n(\mathcal Q_m,A_{\neg S};\ZZ).
\]
\end{corollary}
\begin{proof}
由于 \(C_n(\mathcal Q_m;\ZZ)\) 与 \(\ZZ^{S}\) 均为自由阿贝尔群，\(\ker f_S\) 亦为自由阿贝尔群，且
\(
\rank_{\ZZ}(\ker f_S)=\mathcal L(f_S)
\)
为圆维损失（定义 \ref{def:cdim-audit-completion-loss}）。
结论由定理 \ref{thm:cdim-minimal-audit-completion} 直接给出。证毕。
\end{proof}

\begin{theorem}[部分边界 Gödelization 的最小审计补全与素数审计码]\label{thm:spg-partial-boundary-godel-audit-optimal}
沿用定理 \ref{thm:spg-partial-boundary-screen-kernel-relative-homology} 的记号，并令 \(r:=\rank_{\ZZ}(\ker f_S)\)。
取 \(\ker f_S\) 的一组 \(\ZZ\)-基 \(\{k_1,\dots,k_r\}\)。
再取互异素数 \(\{\pi_1,\dots,\pi_{2r}\}\)，并要求其与用于面标签的素数族 \(\{q_{\vec f}\}\)（定理 \ref{thm:spg-boundary-godelization-holographic-dictionary}）不交。
则存在同态
\[
\mathrm{Enc}_S:\ C_n(\mathcal Q_m;\ZZ)\longrightarrow \ZZ^{S}\times \NN,
\qquad
\mathrm{Enc}_S(x):=\bigl(f_S(x),\ \mathrm{Audit}_S(x)\bigr),
\]
满足：
\begin{enumerate}
  \item \(\mathrm{Enc}_S\) 为单射；
  \item 存在整数向量 \(\alpha(x)=(\alpha_1(x),\dots,\alpha_r(x))\in\ZZ^r\)（关于 \(x\) 的群同态），使得
        \[
        \mathrm{Audit}_S(x)
        :=
        \prod_{j=1}^{r}\pi_j^{\max(0,\alpha_j(x))}\cdot \pi_{j+r}^{\max(0,-\alpha_j(x))},
        \]
        并且由 \(\mathrm{Audit}_S(x)\) 的素因子分解可唯一恢复 \(\alpha(x)\)；
  \item 在所有使 \(f_S\) 可单射提升的审计补全中，所需的独立补账秩（圆维）最小且等于 \(r\)。
\end{enumerate}
\end{theorem}
\begin{proof}
由于 \(C_n(\mathcal Q_m;\ZZ)\) 与 \(\ZZ^{S}\) 为自由阿贝尔群，短正合列
\[
0\longrightarrow \ker f_S\longrightarrow C_n(\mathcal Q_m;\ZZ)\longrightarrow \mathrm{im}(f_S)\longrightarrow 0
\]
以 \(\mathrm{im}(f_S)\) 自由为前提可分裂，从而存在同态重traction \(p:\,C_n(\mathcal Q_m;\ZZ)\to \ker f_S\) 使 \(p|_{\ker f_S}=\mathrm{id}\)。
令 \(\alpha(x)\in\ZZ^r\) 为 \(p(x)\) 在基 \(\{k_j\}\) 下的坐标，即 \(p(x)=\sum_{j=1}^r \alpha_j(x)k_j\)。
则 \(\alpha\) 为群同态，且由算术基本定理 \(\mathrm{Audit}_S(x)\) 的素因子分解唯一确定 \(\alpha(x)\)。
\par
若 \(\mathrm{Enc}_S(x)=\mathrm{Enc}_S(y)\)，则 \(f_S(x-y)=0\) 从而 \(x-y\in\ker f_S\)，并且 \(\alpha(x-y)=\alpha(x)-\alpha(y)=0\)。
由 \(p|_{\ker f_S}=\mathrm{id}\) 得 \(p(x-y)=x-y\)，从而 \(x-y=0\)，故 \(\mathrm{Enc}_S\) 单射。
\par
最优性由推论 \ref{cor:spg-partial-boundary-min-audit-cost-kernel-rank} 给出：任意审计补全的最小成本等于 \(\rank_{\ZZ}(\ker f_S)=r\)，因此不存在秩更小的补全仍能实现单射提升。证毕。
\end{proof}

\begin{theorem}[可见秩函数的可表示拟阵性与审计缺口的超模性]\label{thm:spg-screen-rank-matroid-supermodularity}
沿用上述记号。
对任意 \(S\subset \overrightarrow{\mathcal F_{m}^{n-1}}\)，令
\[
r(S):=\rank_{\QQ}\bigl(\mathrm{im}(f_S\otimes\QQ)\bigr),
\qquad
\delta(S):=\rank_{\QQ}\bigl(\ker(f_S\otimes\QQ)\bigr).
\]
则：
\begin{enumerate}
  \item \(r(\cdot)\) 是 \(\overrightarrow{\mathcal F_{m}^{n-1}}\) 上一个可表示拟阵的秩函数；特别地，对任意 \(S\subset T\) 有 \(r(S)\le r(T)\)，并且对任意 \(S,T\) 有
        \[
        r(S)+r(T)\ \ge\ r(S\cup T)+r(S\cap T).
        \]
  \item \(\delta(S)=\rank_{\QQ}(C_n(\mathcal Q_m;\QQ)) - r(S)\) 从而满足超模性
        \[
        \delta(S)+\delta(T)\ \le\ \delta(S\cup T)+\delta(S\cap T),
        \]
        并且 \(\delta\) 随 \(S\) 单调不增。
  \item 由于 \(C_n(\mathcal Q_m;\ZZ)\) 自由，\(\rank_{\QQ}\ker(f_S\otimes\QQ)=\rank_{\ZZ}\ker f_S\)。
        结合推论 \ref{cor:spg-partial-boundary-min-audit-cost-kernel-rank}，最小审计补全成本 \(S\mapsto \min_{(R,r_S)}\cdim(R)\) 在该屏幕族上同样满足上述超模性。
\end{enumerate}
\end{theorem}
\begin{proof}
取 \(\partial_n\) 在某组 \(\ZZ\)-基下的矩阵表示 \(A\)。
则 \(f_S\otimes\QQ\) 等价于将 \(A\) 的行按 \(S\) 取子集得到的行限制矩阵 \(A_S\)，从而
\[
r(S)=\dim_{\QQ}\Span\{\text{rows of }A\text{ indexed by }S\}.
\]
由线性拟阵的一般理论，\(r(\cdot)\) 满足单调性与次模性，从而为可表示拟阵秩函数。
由秩亏恒等式得 \(\delta(S)=\rank_{\QQ}(C_n(\mathcal Q_m;\QQ))-r(S)\)，因此 \(\delta\) 为超模函数且单调不增。
最后一条由 \(\ker f_S\) 的秩在张量到 \(\QQ\) 后保持不变，以及推论 \ref{cor:spg-partial-boundary-min-audit-cost-kernel-rank} 的同一刻画直接推出。证毕。
\end{proof}

\endinput
